\chapter{Referencia de terminos}

\begin{description}
  \item[Funcion objetivo] Expresion matematica que cuantifica el criterio que
  se desea maximizar o minimizar.
  \item[Variable de decision] Magnitud controlable por el analista o por el
  decisor dentro del modelo.
  \item[Parametro] Dato exogeno que el modelo toma como conocido durante la
  solucion.
  \item[Restriccion] Condicion que limita los valores admisibles de las
  variables de decision.
  \item[Region factible] Conjunto de soluciones que satisfacen todas las
  restricciones.
  \item[Solucion optima] Solucion factible que alcanza el mejor valor de la
  funcion objetivo.
  \item[Relajacion lineal] Modelo obtenido al reemplazar restricciones de
  integralidad por variables continuas.
  \item[Utilizacion] Fraccion de la capacidad de servicio empleada en promedio.
  \item[Makespan] Tiempo de finalizacion del ultimo trabajo en un programa de
  produccion o servicio.
\end{description}
